% ---
% Reemplaza líneas 1093-1138 de Docs/Technical_Report.tex
% Decisiones: N/A (sin cambio)
% ---

\subsubsection{Sistemas de múltiples qubits y entrelazamiento}
% -----------------------------------------------------------

El estado de un sistema de $n$ qubits se describe mediante el producto tensorial de los espacios de Hilbert individuales:

\begin{equation}
    \mathcal{H}^{\otimes n} = 
    \underbrace{\mathbb{C}^2 \otimes \mathbb{C}^2 \otimes 
    \cdots \otimes \mathbb{C}^2}_{n}, 
    \quad \dim(\mathcal{H}^{\otimes n}) = 2^n
    \label{eq:tensor}
\end{equation}

\noindent El estado general de $n$ qubits es una superposición 
de los $2^n$ estados de la base computacional 
$\{|0\cdots0\rangle, |0\cdots1\rangle, \ldots, 
|1\cdots1\rangle\}$:

\begin{equation}
    |\Psi\rangle = \sum_{k=0}^{2^n - 1} c_k |k\rangle, 
    \quad c_k \in \mathbb{C}, 
    \quad \sum_{k=0}^{2^n-1} |c_k|^2 = 1
    \label{eq:multiQubit}
\end{equation}

El \textit{entrelazamiento cuántico} es la propiedad por la cual el estado de un sistema de múltiples qubits no puede factorizarse como el producto tensorial de estados individuales. Un estado entrelazado no tiene representación separable y exhibe correlaciones entre qubits que no tienen análogo clásico \cite{nielsen2010quantum}. La compuerta estándar para generar entrelazamiento es la \textit{compuerta CNOT} (Controlled-NOT), que opera sobre dos qubits: un qubit de control $c$ y un qubit objetivo $t$, aplicando una negación al objetivo condicionada al estado del control:

\begin{equation}
    \text{CNOT} = 
    \begin{pmatrix} 
        1 & 0 & 0 & 0 \\ 
        0 & 1 & 0 & 0 \\ 
        0 & 0 & 0 & 1 \\ 
        0 & 0 & 1 & 0 
    \end{pmatrix}
    \label{eq:cnot}
\end{equation}

\noindent La combinación de compuertas de rotación de un qubit 
y compuertas CNOT es suficiente para aproximar cualquier 
operación unitaria sobre $n$ qubits, propiedad conocida como 
\textit{universalidad} \cite{nielsen2010quantum}. Esta 
combinación es precisamente la que estructura el ansatz 
variacional RealAmplitudes utilizado en el presente trabajo.
